\documentclass[12pt]{article}

\usepackage[letterpaper,margin=1in]{geometry}
\usepackage[T1]{fontenc}
\usepackage[utf8]{inputenc}
\usepackage{amsmath}
\usepackage{array}
\usepackage{booktabs}
\usepackage{caption}
\usepackage{graphicx}
\usepackage{hyperref}

\hypersetup{
  colorlinks=true,
  linkcolor=black,
  citecolor=black,
  urlcolor=blue
}

\newcommand{\tables}{../Output/Tables}
\newcommand{\figures}{../Output/Figures}

% Restore captions from Memo/Rozee_AI.pdf while still using the table source files.
\newcommand{\inputtablewithcaption}[2]{%
  \begingroup
  \let\savedcaption\caption
  \renewcommand{\caption}[1]{\savedcaption{#1}}%
  \input{#2}%
  \endgroup
}

\title{Rozee: AI \& Pakistan Labor Market}
\author{Yifei Wu, Converted to Latex by HW}
\date{Feb 2026}

\begin{document}

\maketitle
\tableofcontents
\listoffigures
\listoftables
\clearpage

\section{Introduction}

This project utilizes data from Rozee, Pakistan's leading recruitment platform, to
investigate the impact of AI on the labor market. The dataset provides
comprehensive coverage, encompassing detailed information on users, companies, job
postings, and job applications.

At this stage, we have constructed the AI exposure index at the job-posting level,
encompassing around 300k postings from 2020 through April 2025; for details on
the methodology, please refer to \href{https://www.dropbox.com/scl/fi/uqi6zjdkaz9e7wlg5p04j/AI-Pakistan-LM-Presentation.pptx?rlkey=2xudrhus8cz26kf3xi7oix7fp&e=2&st=mdsvydvb&dl=0}{here}. Building on this, we will proceed with a
series of descriptive statistics and empirical analyses.

\section{Descriptives}

\subsection*{AI Exposure}

Please see Table \ref{tab:exposure_categories}, \ref{tab:job_level_exposure},
\ref{tab:ai_exposure_by_sector}; Figure \ref{fig:ai_exposure_cdf},
\ref{fig:ai_exposure_characteristics}, \ref{fig:ai_exposure_tasks}.

\begin{table}[h]
\centering
\caption{Exposure Categories for Each Task}
\label{tab:exposure_categories}
\begin{tabular}{lll}
\toprule
Category & Rubric & Explanation \\
\midrule
E0 & No Exposure & Minimal reduction in time or decreases quality \\
E1 & Direct Exposure & Access to LLM alone or simple interface decreases time \\
   &                 & to complete task by 50\% \\
E2 & LLM + Exposed  & Additional software needed on top of LLM \\
\bottomrule
\end{tabular}
\end{table}

\begin{table}[h]
\centering
\caption{Job-Level Exposure Measures}
\label{tab:job_level_exposure}
\begin{tabular}{ll}
\toprule
Variable & Definition \\
\midrule
$\alpha$ & Share of job's tasks labelled E1 \\
$\beta$  & Share of job's tasks labelled E1 + 0.5$\times$ Share of job's tasks labelled E2 \\
$\gamma$ & Share of job's tasks labelled E1 + Share of job's tasks labelled E2 \\
\bottomrule
\end{tabular}
\end{table}

\section{Analysis}

Using job-level AI exposure data, we investigate the broad effects of AI on the
labor market. Now we consider the full-scale rollout of ChatGPT in November 2022
as a key exogenous shock for our analysis.

\begin{figure}[p]
\centering
\caption{AI-exposed tasks in Job Postings ($\beta$)}
\label{fig:ai_exposure_cdf}
\includegraphics[width=0.85\textwidth]{\figures/ai_exposure_cdf.png}
\end{figure}

\begin{table}[p]
\centering
\caption{AI exposure by ISIC sector ($\beta$)}
\label{tab:ai_exposure_by_sector}
\small
% Source: \tables/ai_exposure_by_sector.xlsx
\resizebox{\textwidth}{!}{%
\begin{tabular}{lrrrrr}
\toprule
ISIC Sector & \begin{tabular}{c}Mean\\(Human)\end{tabular} &
\begin{tabular}{c}SD\\(Human)\end{tabular} &
\begin{tabular}{c}Mean\\(GPT)\end{tabular} &
\begin{tabular}{c}SD\\(GPT)\end{tabular} & N \\
\midrule
Accommodation and food service activities & 0.50 & 0.21 & 0.52 & 0.21 & 5,082 \\
Administrative and support service activities & 0.54 & 0.17 & 0.61 & 0.18 & 60,359 \\
Agriculture, forestry and fishing & 0.48 & 0.19 & 0.55 & 0.19 & 841 \\
Arts, entertainment and recreation & 0.51 & 0.19 & 0.56 & 0.20 & 2,477 \\
Construction & 0.47 & 0.17 & 0.53 & 0.17 & 4,387 \\
Education & 0.51 & 0.17 & 0.54 & 0.19 & 15,694 \\
Electricity, gas, steam and air conditioning supply & 0.46 & 0.21 & 0.51 & 0.19 & 1,304 \\
Financial and insurance activities & 0.54 & 0.14 & 0.59 & 0.13 & 13,461 \\
Human health and social work activities & 0.46 & 0.18 & 0.51 & 0.17 & 24,586 \\
Information and communication & 0.57 & 0.14 & 0.72 & 0.19 & 93,994 \\
Manufacturing & 0.47 & 0.18 & 0.54 & 0.18 & 40,615 \\
Mining and quarrying & 0.48 & 0.16 & 0.54 & 0.17 & 1,125 \\
Other service activities & 0.46 & 0.22 & 0.50 & 0.21 & 561 \\
Professional, scientific and technical activities & 0.54 & 0.17 & 0.60 & 0.18 & 32,953 \\
Public administration and defence & 0.50 & 0.18 & 0.56 & 0.20 & 560 \\
Real estate activities & 0.50 & 0.17 & 0.56 & 0.15 & 13,419 \\
Transportation and storage & 0.49 & 0.17 & 0.58 & 0.17 & 4,683 \\
Wholesale and retail trade & 0.52 & 0.17 & 0.58 & 0.17 & 11,718 \\
\midrule
\textbf{Total (Job Postings)} & \textbf{0.53} & \textbf{0.17} & \textbf{0.61} & \textbf{0.19} & \textbf{327,819} \\
\bottomrule
\end{tabular}%
}
\end{table}

\subsection{Labor Demand}

Here we focus on the impact of AI on the demand side of the labor market.
Specifically, we look at changes in how job postings are designed and how
companies manage their posting behaviors.

\subsection*{Job Counts and Wages}

We aim to investigate whether job posting volumes and wages exhibit significant
trends following the AI shock. Building on existing methodologies, we aggregate
individual job postings into different O*NET occupation levels. This allows us
to construct an occupation-level AI exposure index, along with the total count of
job postings and average wages at the occupation-month level. Following
Teutloff et al. (2025), we then estimate the following continuous DiD regression:

\[
Y_{it}=\beta_0+\beta_1(LLM\_Exposure_i \times Post_t)+\gamma_i+\theta_t+\varepsilon_{it}
\]

As shown in Tables \ref{tab:did_postings_gpt}, \ref{tab:did_postings_human},
\ref{tab:did_wages_gpt}, \ref{tab:did_wages_human}, occupations with high AI
exposure experienced a contraction in job posting volumes and lower wages
following the shock compared to low-exposure occupations. The event study in
Figure \ref{fig:es_postings} and Figure \ref{fig:es_wages} further confirms the
parallel trends assumption and validates our primary findings.

\begin{figure}[p]
\centering
\caption{AI exposure by Job Salary/Education/Experience ($\beta$)}
\label{fig:ai_exposure_characteristics}
\includegraphics[width=0.55\textwidth]{\figures/ai_exposure_salary.png}

\vspace{0.5em}
{\large Job Posting Salary}

\vspace{1.5em}
\begin{minipage}{0.47\textwidth}
\centering
\includegraphics[width=\textwidth]{\figures/ai_exposure_education.png}
\vspace{0.5em}

{\large Education Requirement}
\end{minipage}\hfill
\begin{minipage}{0.47\textwidth}
\centering
\includegraphics[width=\textwidth]{\figures/ai_exposure_experience.png}
\vspace{0.5em}

{\large Experience Requirement}
\end{minipage}
\end{figure}

\begin{figure}[p]
\centering
\caption{AI exposure measure \& Number of tasks extracted ($\beta$)}
\label{fig:ai_exposure_tasks}
\begin{minipage}{0.47\textwidth}
\centering
\includegraphics[width=\textwidth]{\figures/ai_exposure_num_tasks_gpt.png}
\end{minipage}\hfill
\begin{minipage}{0.47\textwidth}
\centering
\includegraphics[width=\textwidth]{\figures/ai_exposure_num_tasks.png}
\end{minipage}
\end{figure}

\inputtablewithcaption{Impact of public-use ChatGPT on job counts (GPT rating)\label{tab:did_postings_gpt}}{\tables/did_postings_gpt.tex}

\inputtablewithcaption{Impact of public-use ChatGPT on job counts (Human rating)\label{tab:did_postings_human}}{\tables/did_postings_human.tex}

\begin{figure}[p]
\centering
\caption{Event study of AI exposure on Job Counts (Human rating)}
\label{fig:es_postings}
\includegraphics[width=0.75\textwidth]{\figures/es_postings.png}
\end{figure}

\begin{figure}[p]
\centering
\caption{Event study of AI exposure on Wages (Human rating)}
\label{fig:es_wages}
\includegraphics[width=0.75\textwidth]{\figures/es_wages.png}
\end{figure}

\inputtablewithcaption{Impact of public-use ChatGPT on wages (GPT rating)\label{tab:did_wages_gpt}}{\tables/did_wages_gpt.tex}

\inputtablewithcaption{Impact of public-use ChatGPT on wages (Human rating)\label{tab:did_wages_human}}{\tables/did_wages_human.tex}

\begin{thebibliography}{9}

\bibitem{teutloff2025}
Teutloff, O., Einsiedler, J., K\"assi, O., Braesemann, F., Mishkin, P., and
del Rio-Chanona, R. M. (2025). Winners and losers of generative ai: Early
evidence of shifts in freelancer demand. \emph{Journal of Economic Behavior \&
Organization}, page 106845.

\end{thebibliography}

\end{document}
